\section*{Capítulo \ref{ch:g12:photosynthesis} --- La fotosíntesis}

\begin{solution}{exo:g12:photosynthesis:1}
Un \omterm{def:g10:cells-common-unit:organelle}{orgánulo} de doble membrana con un estroma líquido que contiene pilas
de membranas tilacoidales. Los pigmentos y las \omterm{def:g10:chemistry-of-life:families}{proteínas} que convierten
la luz están en las membranas de los \omterm{def:g12:photosynthesis:chloroplast}{tilacoides}; las \omterm{def:g11:enzymes-and-phenotype:enzyme}{enzimas} que
construyen azúcar, en el estroma.
\end{solution}

\begin{solution}{exo:g12:photosynthesis:2}
La clorofila absorbe el azul y el rojo y refleja el verde, que es lo
que vemos. La luz azul y la roja son las que mejor impulsan la
\omterm{def:g10:cell-metabolism:autotrophy}{fotosíntesis}.
\end{solution}

\begin{solution}{exo:g12:photosynthesis:3}
ATP, NADP reducido y oxígeno; el oxígeno es el subproducto, que se
libera.
\end{solution}

\begin{solution}{exo:g12:photosynthesis:4}
Del agua que se rompe, no del dióxido de carbono: unas algas a las que
se dio agua marcada con oxígeno pesado liberaron $\mathrm{O_2}$ pesado;
con dióxido de carbono marcado, no.
\end{solution}

\begin{solution}{exo:g12:photosynthesis:5}
Seis $\mathrm{CO_2}$, 18 ATP y 12 NADP reducido.
\end{solution}

\begin{solution}{exo:g12:photosynthesis:6}
A \qty{550}{nm}: absorción de un 8\%, velocidad de un 25\%. A
\qty{665}{nm}: absorción de un 90\%, velocidad de un 95\%. Las lámparas
verdes entregan la luz que la hoja menos usa; el invernadero ahorraría
energía sin cultivar casi nada.
\end{solution}

\begin{solution}{exo:g12:photosynthesis:7}
Las bacterias nadan hacia el oxígeno; se juntan donde se está liberando
oxígeno, es decir, donde el alga hace la \omterm{def:g10:cell-metabolism:autotrophy}{fotosíntesis}, y eso es bajo
las partes azul y roja del espectro. La \omterm{def:g10:cell-metabolism:autotrophy}{fotosíntesis}, medida por su
oxígeno, depende del color de la luz y sigue la absorción de los
pigmentos.
\end{solution}

\begin{solution}{exo:g12:photosynthesis:8}
Se libera oxígeno sin que se fije carbono alguno: la fase luminosa, que
rompe el agua y mueve electrones, está separada de la fase del carbono
y no necesita $\mathrm{CO_2}$; necesita luz, agua y algo que reciba los
electrones.
\end{solution}

\begin{solution}{exo:g12:photosynthesis:9}
Sin luz no hay ATP ni NADP reducido, así que el ácido de tres carbonos
no puede reducirse y se acumula; el aceptor sigue capturando
$\mathrm{CO_2}$ durante un rato, pero ya no se regenera, y se agota.
\end{solution}

\begin{solution}{exo:g12:photosynthesis:10}
La intensidad de la luz baja con la distancia (más o menos con su
cuadrado: cuatro veces menos al doble de distancia), y la fase luminosa
va más despacio. A oscuras no se libera oxígeno --- y la planta consume
algo por respiración.
\end{solution}

\begin{solution}{exo:g12:photosynthesis:11}
$8/44 \approx \qty{0.18}{mol}$ de $\mathrm{CO_2}$: $0.18/6 \approx
\qty{0.030}{mol}$ de glucosa, unos \qty{5.5}{g}; \qty{0.18}{mol} de
$\mathrm{O_2}$, unos \qty{5.8}{g}.
\end{solution}

\begin{solution}{exo:g12:photosynthesis:12}
Ella misma no usa luz, solo los productos de la fase luminosa. Pero en
una planta que se deja a oscuras se detiene en un minuto, en cuanto se
gastan el ATP y el NADP reducido: funciona con luz, alimentada por la
fase luminosa, y ``oscura'' describe lo que necesita, no cuándo ocurre.
\end{solution}

\begin{solution}{exo:g12:photosynthesis:13}
$^{18}$O: en el $\mathrm{O_2}$ liberado al cabo de un minuto, y ahí
sigue al cabo de una hora y de un día (algo también en el agua que
produce la respiración). $^{14}$C: al cabo de un minuto, en el ácido de
tres carbonos y en los azúcares del ciclo; al cabo de una hora, en
sacarosa y en almidón; al cabo de un día, en celulosa, \omterm{def:g10:chemistry-of-life:families}{lípidos},
\omterm{def:g10:chemistry-of-life:families}{aminoácidos} y en el $\mathrm{CO_2}$ expulsado por la respiración.
\end{solution}

\begin{solution}{exo:g12:photosynthesis:14}
$20/1000 = 2\%$. El resto se refleja (el verde), se absorbe en las
longitudes de onda equivocadas y se convierte en calor, se gasta en la
transpiración que evapora agua, o se pierde como calor en las propias
reacciones luminosas.
\end{solution}

\begin{solution}{exo:g12:photosynthesis:15}
La respiración, la combustión y la oxidación de las rocas y del hierro
consumen oxígeno continuamente; sin una \omterm{def:g10:cell-metabolism:autotrophy}{fotosíntesis} que lo renueve, la
reserva de la atmósfera se agotaría en unos pocos miles de años. La
atmósfera primitiva no tenía oxígeno libre; solo se acumuló después de
que los organismos fotosintéticos llevaran muchísimo tiempo
produciéndolo --- los primeros dos mil millones de años fueron pobres
en oxígeno.
\end{solution}

\begin{solution}{pb:g12:photosynthesis:1}
\textbf{1.} $8/44 \approx \qty{0.18}{mol}$.

\textbf{2.} $0.18/6 = \qty{0.030}{mol}$, unos \qty{5.5}{g}.

\textbf{3.} \qty{0.18}{mol} de $\mathrm{O_2}$, unos \qty{4.4}{L}.

\textbf{4.} Dos moléculas de agua por $\mathrm{O_2}$: \qty{0.36}{mol},
unos \qty{6.5}{g} --- un uno por ciento de los \qty{600}{mL}
transpirados.

\textbf{5.} De las moléculas de agua rotas en los \omterm{def:g12:photosynthesis:chloroplast}{tilacoides}: lo
muestra el marcado con oxígeno pesado, que aparece en el
$\mathrm{O_2}$ solo cuando el agua está marcada.

\textbf{6.} Una vuelta por $\mathrm{CO_2}$: \qty{0.18}{mol} de vueltas.

\textbf{7.} Tres ATP y dos NADP reducido por vuelta: \qty{0.55}{mol} de
ATP y \qty{0.36}{mol} de NADP reducido.

\textbf{8.} Cada agua rota da dos electrones, o sea un NADP reducido:
\qty{0.36}{mol} de agua --- la misma cifra que en la pregunta 4, como
tenía que ser.

\textbf{9.} El oxígeno se detendría de inmediato; el NADP reducido y el
ATP dejarían de fabricarse y se agotarían en unos segundos; el azúcar
se detendría en cuanto se acabaran los transportadores.

\textbf{10.} El \omterm{prop:g12:photosynthesis:calvin}{ciclo de Calvin} ya no gastaría los transportadores; sin
nada que recibiera sus electrones y su energía, la fase luminosa se
embotellaría y se frenaría en unos segundos --- las dos fases están
acopladas por los transportadores.

\textbf{11.} $0.030 \times 2870 \approx \qty{86}{kJ}$ por hora, es
decir, $86000/3600 \approx \qty{24}{W}$.

\textbf{12.} $24/1000 = 2.4\%$ de la luz solar total; $24/450 \approx
5\%$ de las longitudes de onda absorbidas.

\textbf{13.} $8 \times 176 \approx \qty{1400}{kJ}$ de fotones por mol
de $\mathrm{CO_2}$; almacenados: $2870/6 \approx \qty{480}{kJ}$.

\textbf{14.} Alrededor de un tercio; el resto se libera como calor en
las transferencias sucesivas de energía de las dos fases.

\textbf{15.} Ganancia neta: $5.5 - 0.5 = \qty{5}{g}$ de glucosa por
hora. Se mide como la absorción neta de $\mathrm{CO_2}$ con luz; la
respiración sola se mide a oscuras y se vuelve a sumar para obtener la
\omterm{def:g10:cell-metabolism:autotrophy}{fotosíntesis} bruta.

\textbf{16.} $\num{8.5e11}/\num{1e11} \approx 8.5$ años --- y, con las
algas del océano fijando otro tanto, unos 4 años.

\textbf{17.} La respiración (de plantas, animales, hongos y bacterias),
junto con el fuego, devuelve dióxido de carbono casi al mismo ritmo;
fijación y devolución se equilibraban, así que la reserva era estable
hasta que la quema de combustibles fósiles añadió una fuente nueva.

\textbf{18.} \qty{1e11}{t} de carbono son $\num{8.3e12}\,\unit{mol}$;
ese mismo número de moles de $\mathrm{O_2}$ son unas \qty{2.7e11}{t} de
oxígeno.

\textbf{19.} $\num{1.2e15}/\num{2.7e11} \approx 4500$ años de
producción (unos 2000 con la parte del océano). Pero la respiración
consume oxígeno casi al mismo ritmo al que la \omterm{def:g10:cell-metabolism:autotrophy}{fotosíntesis} lo fabrica,
de modo que la reserva es un equilibrio y no una acumulación: el
oxígeno del aire se formó a lo largo de cientos de millones de años por
el pequeño exceso de la producción sobre el consumo.

\textbf{20.} Alrededor del 2\% de la luz solar total (el 5\% de las
longitudes de onda aprovechables) se almacena como azúcar; el oxígeno
liberado es el oxígeno del agua rota; más o menos una cuarta parte del
carbono de la atmósfera pasa por la \omterm{def:g10:cell-metabolism:autotrophy}{fotosíntesis} cada año.
\end{solution}
